\begin{frame}{데이터 고르기 실전 가이드 5개}
  후보에 ``점수''를 매겨보면 팀 회의가 쉬워진다:
  \begin{enumerate}
    \item \kb{행의 물리적 단위}를 한 문장으로 설명할 수 있는가
      (환자? 거래? 주문?) --- 설명이 안 되면 ``예측 대상 단위를
      일반화하는가''를 물을 수 없다.
    \item \kb{레이아웃의 함정}이 있는가 --- 같은 환자의 여러 검사
      결과가 여러 행이면 무작위 분할이 누수 $\to$
      \texttt{GroupKFold} 필요.
    \item \kb{시계열은 시간순}으로 나눠야 한다(6.3절 노트북 §5).
    \item \kb{클래스 비율을 먼저 보고 지표를 정한다} --- 99:1이면
      정확도가 아니라 PR-AUC/재현율(Week 2.3 정확도의 함정).
    \item \kb{크기} --- 100행 미만이면 CV 폴드가 너무 작아 숫자가
      극도로 흔들림 $\to$ 500$\sim$10{,}000행이 편함.
  \end{enumerate}
\end{frame}
